\chapter{Discussion and Conclusion}

\section{Interpretation of Results}

[To be completed: high-level reading of the experimental findings, unexpected patterns, and what the results suggest about the relationship between ambiguity and classification performance.]

\section{Answers to Research Questions}

\subsection{RQ1: Does ambiguity negatively affect BERT classification performance?}

[To be completed.]

\subsection{RQ2: Can LLM disambiguation improve BERT classification accuracy?}

[To be completed.]

\subsection{RQ3: Can LLMs match or exceed BERT as direct classifiers?}

RQ3 evaluates whether a general-purpose LLM can serve as a standalone classifier, replacing the fine-tuned BERT pipeline entirely. Three GPT-family models of varying capability (GPT\textsubscript{4.1-nano}, GPT\textsubscript{5.4-mini}, GPT\textsubscript{5.5}) are tested in both zero-shot and few-shot configurations on the full 2,166-sentence dataset. Performance is compared against the BERT baseline using the same multi-label metrics (exact-match accuracy, Hamming loss, Jaccard score, micro/macro/per-label F1) and binary requirement F1-score.

[To be completed: results and interpretation.]

\section{Threats to Validity and Limitations}

[To be completed: dataset size, single-annotator labeling, LLM non-determinism, generalisability beyond cloud computing domain.]

\section{Future Work}

[To be completed: multi-annotator studies, other domains, fine-tuned LLM classifiers, automated ambiguity detection.]

\section{Conclusion}

[To be completed: brief summary of the thesis contributions and closing remarks.]
